
\uline{on the other hand}, (IV 2.11) \uline{one has a homomorphism}
@@M0@@
\uline{Then}, \uline{the composite}
@@M1@@
\uline{is zero}.

Indeed, @@M2@@ is a regular immersion of codimension @@M3@@ whose conormal sheaf @@M4@@ is \emph{free} of rank @@M5@@, since @@M6@@ can be described by a global equation @@M7@@ (where @@M8@@ is a uniformizer of @@M9@@). Thus one has @@M10@@, and the conclusion follows from 7.1.

Throughout what follows, we shall suppose that @@M11@@ is \emph{noetherian} and \emph{flat over} @@M12@@. Let us now apply the exact homotopy sequence (IX 1.1) to the following situation: @@M13@@, with the open subscheme @@M14@@ and complementary closed subscheme @@M15@@:
@@M16@@
where @@M17@@ is the inclusion, which is an open immersion. Using 7.2, one obtains:

\sgaStableTarget{sga6.e10.prop.7-3}\textbf{Proposition 7.3.} \uline{With the notation of} 7.2, @@M18@@ \uline{being assumed noetherian}, \uline{there is a unique group homomorphism}
\sgaStableTarget{sga6.e10.eq.7-3-1}@@M19@@
(\uline{called the} ``specialization homomorphism'') \uline{making the diagram commutative}
\sgaStableTarget{sga6.e10.eq.7-3-2}@@M20@@


% Local macro additions for Exposé X continuation
\providecommand{\Kalg}{K_{\mathrm{alg}}}
\providecommand{\Gralg}{\operatorname{Gr}_{\mathrm{alg}}}
\providecommand{\Grnum}{\operatorname{Gr}_{\mathrm{num}}}
\providecommand{\Knum}{K_{\mathrm{num}}}
\providecommand{\Cl}{\operatorname{Cl}}
\providecommand{\Ch}{\operatorname{Ch}}
\providecommand{\Nm}{\operatorname{N}}
\providecommand{\colim}{\varinjlim}
\providecommand{\Hens}{\operatorname{Hens}}
\providecommand{\Spec}{\operatorname{Spec}}
\providecommand{\coker}{\operatorname{coker}}

